% ============================================================================
% campos-formulario.tex
% ----------------------------------------------------------------------------
% Define COMANDOS PRONTOS para criar campos de formulário PDF interativos
% (caixas de texto, checkboxes) dentro do documento, usando o pacote
% "hyperref" (que fornece \TextField e \CheckBox).
%
% Esses campos SÓ FUNCIONAM COMO INTERATIVOS quando o PDF é aberto em um
% leitor compatível com formulários (ex.: Adobe Acrobat Reader). Em alguns
% visualizadores simples (como o do navegador), os campos podem aparecer
% apenas como retângulos vazios, sem interatividade.
%
% COMANDOS DISPONÍVEIS PARA USAR NAS SEÇÕES DE CONTEÚDO:
%   \campotexto{nome}{largura}          -> campo de texto de uma linha
%   \campolongo{nome}{altura}           -> campo de texto de múltiplas linhas
%   \caixa{nome}{rótulo}                -> checkbox com texto ao lado
%   \campotabela[largura]{nome}{altura} -> campo de texto para uso dentro de tabelas
%   \checktabela{nome}                  -> checkbox para uso dentro de tabelas
%
% Exemplo de uso em uma seção de conteúdo:
%   Nome: \campotexto{nome_completo}{6cm}
%   \caixa{aceite_termos}{Li e aceito os termos}
% ============================================================================

% -- Aparência padrão dos campos -----------------------------------------------
\definecolor{fieldbg}{RGB}{245, 245, 245}   % cor de fundo cinza-clara dos campos
\newcommand{\tamanhoFonteCampo}{12pt}        % tamanho de fonte dos campos normais
\newcommand{\tamanhoFonteCampoTabela}{10pt}  % tamanho de fonte dos campos dentro de tabelas
\newcommand{\ladoCaixa}{12pt}                % tamanho (lado) dos checkboxes

% Necessário para o correto funcionamento do pacote hyperref com campos de
% formulário; não altere estas duas linhas.
\def\LayoutTextField#1#2{#2}
\def\LayoutChoiceField#1#2{#2}

% -- \campotexto{nome}{largura} ------------------------------------------------
% Cria um campo de texto de uma única linha.
%   #1 = nome interno do campo (identificador único, sem espaços)
%   #2 = largura do campo (ex.: 5cm)
\newcommand{\campotexto}[2]{
    \rule{0pt}{16pt}  % garante altura mínima de linha para alinhamento visual
    \TextField[name=#1, width=#2, height=14pt, borderwidth=0.75,
        bordercolor={0.8 0.8 0.8}, backgroundcolor=fieldbg,
        charsize=\tamanhoFonteCampo]{}
}

% -- \campolongo{nome}{altura} -------------------------------------------------
% Cria um campo de texto de múltiplas linhas (ocupa a largura toda da página).
%   #1 = nome interno do campo
%   #2 = altura do campo (ex.: 3cm)
\newcommand{\campolongo}[2]{
    \noindent\rule{0pt}{#2}
    \TextField[name=#1, width=\linewidth, height=#2, multiline=true,
        borderwidth=0.75, bordercolor={0.8 0.8 0.8}, backgroundcolor=fieldbg,
        charsize=\tamanhoFonteCampo]{}
}

% -- \caixa{nome}{rótulo} -------------------------------------------------------
% Cria um checkbox com um texto/rótulo ao lado.
%   #1 = nome interno do campo
%   #2 = texto exibido ao lado do checkbox
\newcommand{\caixa}[2]{
    {\fontsize{\tamanhoFonteCampo}{\tamanhoFonteCampo}\selectfont
    \CheckBox[name=#1, bordercolor={0 0 0},
        checkboxsymbol=\ding{51},
        width=\ladoCaixa, height=\ladoCaixa]{ #2}}
}

% -- \campotabela[largura]{nome}{altura} ---------------------------------------
% Versão do campo de texto pensada para uso DENTRO DE CÉLULAS DE TABELA.
% O parâmetro de largura é opcional (padrão: \hsize, ou seja, a largura
% disponível da célula).
%   [#1] = largura (opcional, padrão \hsize)
%   #2   = nome interno do campo
%   #3   = altura do campo
\newcommand{\campotabela}[3][\hsize]{
    \makebox[\hsize][c]{
        \TextField[name=#2, width=#1, height=#3, multiline=true,
            borderwidth=0, backgroundcolor=fieldbg,
            charsize=\tamanhoFonteCampoTabela]{}
    }
}

% -- \checktabela{nome} ---------------------------------------------------------
% Checkbox pensado para uso dentro de células de tabela (sem rótulo ao lado,
% e sem borda visível).
%   #1 = nome interno do campo
\newcommand{\checktabela}[1]{
    \makebox[\hsize][c]{
        {\fontsize{\tamanhoFonteCampoTabela}{\tamanhoFonteCampoTabela}\selectfont
        \CheckBox[name=#1, checkboxsymbol=\ding{51},
            width=\ladoCaixa, height=\ladoCaixa, borderwidth=0]{}}
    }
}